\usepackage[utf8]{inputenc} % Allows the use of some special characters
\usepackage{amsmath, amsthm, amssymb, amsfonts} % Nicer mathematical typesetting
% \usepackage{lipsum} % Creates dummy text lorem ipsum to showcase typsetting 

\usepackage{graphicx} % Allows the use of \begin{figure} and \includegraphics
\usepackage{float} % Useful for specifying the location of a figure ([H] for ex.)
\usepackage{caption} % Adds additional customization for (figure) captions
\usepackage{subcaption} % Needed to create sub-figures
% \usepackage[most]{tcolorbox}
\usepackage{tabularray}
% \usepackage{tabularx} % Adds additional customization for tables
% \usepackage{tabu} % Adds additional customization for tables
\usepackage{booktabs} % For generally nicer looking tables
\usepackage{longtable} %longtable
\usepackage{multirow}
% \usepackage[nottoc,numbib]{tocbibind} % Automatically adds bibliography to ToC
\usepackage[margin = 2.5cm]{geometry} % Allows for custom (wider) margins
% \usepackage{microtype} % Slightly loosens margin restrictions for nicer spacing  
\usepackage{titlesec} % Used to create custom section and subsection titles
\usepackage{titletoc} % Used to create a custom ToC
\usepackage{pifont} %For extra symbol
\usepackage{fancyhdr} % Adds customization for headers and footers
\usepackage[shortlabels]{enumitem} % Adds additional customization for itemize. 

\usepackage{pdfpages} %% for pdf
\usepackage{pgfplots}%for pdf plot
\pgfplotsset{width=10cm,compat=1.9}
%\usepackage{circuitikz}


\usepackage{hyperref} % Allows links and makes references, and the ToC clickable
\usepackage[noabbrev, capitalise]{cleveref} % Easier referencing using \cref{<label>} instead of \ref{}

\usepackage{xcolor} % Predefines additional colors and allows user-defined colors

\usepackage{adjustbox} % Fitting figures and tables inside the page width

\usepackage{mathtools} % To use box inside an aligned environment

% Allotting space for the header
\setlength{\headheight}{25pt}%

% Defines a command used by tikz to calculate some coordinates for the front-page
% \makeatletter
% \newcommand{\gettikzxy}[3]{%
%   \tikz@scan@one@point\pgfutil@firstofone#1\relax
%   \edef#2{\the\pgf@x}%
%   \edef#3{\the\pgf@y}%
% }
% \makeatother

% \usepackage{tikz} % Useful for drawing images, used for creating the frontpage
% \usetikzlibrary{positioning} % Additional library for relative positioning 
% \usetikzlibrary{calc} % Additional library for calculating within tikz
% \usetikzlibrary{shapes.geometric, arrows}

% Defining Tickz Style
% \tikzstyle{startstop} = [rectangle, rounded corners, minimum width=2cm, minimum height=1cm, text centered, text width = 8cm, draw=black, fill=white]

% \tikzstyle{io} = [trapezium, trapezium left angle=70, trapezium right angle=110, minimum width=3cm, minimum height=1cm, text centered, text width = 4.5cm, draw=black, fill=blue!30]

% \tikzstyle{process} = [rectangle, minimum width=2cm, minimum height=1cm, text centered, text width = 8cm, draw=black, fill=white]

% \tikzstyle{decision} = [diamond, minimum width=3cm, minimum height=1cm, text centered, draw=black, fill=green!30]

% \tikzstyle{arrow} = [ultra thick,->,>=stealth]

\newcommand{\name}[1]{#1\\}
\NewDocumentEnvironment{distlist}{+b}{%
  \begin{center}
  \begin{tblr}[expand=\name]{
    width=0.8\textwidth,
    colspec={X[1,l] X[1,l]},
    hline{1} = {2pt},
    hline{2} = {2pt},
    hline{3-Z} = {0.5pt,dotted},
    vline{2} = {dotted},
    % row{1} = {1.1cm},
  }
  Name & Roll Number \\
  #1
  \end{tblr}
  \end{center}
}{}

